\section{Excluding conjugacy in minimizing convex hulls}
\label{sec:conjugacy}

The continuation argument will encounter finite convex hulls whose every
velocity is known to be minimizing. Their generators will be nonconjugate,
but their other points may lie at cut. We show here that none of those
points can be conjugate. This removes the focal obstruction before the
next section addresses uniqueness of minimizing geodesics.

The essential step concerns a segment with nonconjugate endpoints.
Transverse concavity first confines a
hypothetical conjugacy kernel to the segment direction. The radial
Hessian identity then controls that remaining direction and contradicts
the singular growth forced by an index-form trial field. All uses of
weak MTW occur at proper radial contractions, where the true cost is
smooth.

\begin{proposition}
\label{prop:segment-nonconjugacy}
Suppose that $[b_0,b_1]\subset G_x$ and that $b_0,b_1$ are
nonconjugate.  Then every velocity in $[b_0,b_1]$ is nonconjugate.
\end{proposition}

\begin{proof}
The assertion is immediate if $b_0=b_1$.  Otherwise, suppose there
is a singular interior velocity $z$, and write the segment as
\[
 z+te,\qquad -l_0\leq t\leq l_1,
 \qquad |e|_x=1,\quad l_0,l_1>0.
\]
Here $b_0=z-l_0e$ and $b_1=z+l_1e$.  The velocity $z$ is nonzero,
since $d_0\exp_x=\Id$.
For every $0<r<1$, the full contracted segment
$\{r(z+te):-l_0\leq t\leq l_1\}$ is contained in $I(x)$ by
\eqref{eq:radial-interior}.

Choose $0\ne k\in\ker d_z\exp_x$.  Variation of the initial
velocity in the direction $k$ gives a Jacobi field $J_z$ along
$\gamma_{x,z}$ on $[0,1]$ with
\[
 J_z(0)=J_z(1)=0,\qquad D_\tau J_z(0)=k.
\]
The function $\langle J_z,\dot\gamma_{x,z}\rangle$ is affine in
$\tau$ and vanishes at both endpoints.  Its initial derivative is
$\langle k,z\rangle_x$, so $k\perp z$.

\medskip\noindent
\emph{A trial-field estimate.}
The Jacobi-field trial below is related to the conjugate-cut argument of
Cordero-Erausquin, McCann and Schmuckenschl\"ager
\cite[Proposition~2.5]{CMS2001}. We derive the needed Hessian estimate
explicitly and combine it with transverse MTW concavity and the radial
identity.

Fix a vector $\xi\in T_xM$ such that $\langle\xi,k\rangle_x\ne0$.
For a velocity $a$ near $z$, let $T_a(\tau)$ denote parallel
transport from $x$ along $\gamma_{x,a}$ to time $\tau$, and define
\[
 j(\tau)=T_z(\tau)^{-1}J_z(\tau),\qquad
 J_a(\tau)=T_a(\tau)j(\tau),\qquad
 E_a(\tau)=(1-\tau)T_a(\tau)\xi.
\]
The geodesic and parallel-transport equations give smooth dependence
on $(a,\tau)$.  Thus these are smooth families of fields along the
parametrized geodesics, with
$J_a(0)=J_a(1)=0$ and $E_a(0)=\xi$, $E_a(1)=0$.
Every $J_a$ is nonzero, since its components $j(\tau)$ are fixed and
not identically zero.  The construction uses no inverse of the
exponential map at $z$.

Write $Q_a=Q_{\gamma_{x,a}}$ on $[0,1]$ and set
$P(a)=Q_a(J_a,J_a)$.  This is a smooth function on a full
neighborhood of $z$, and $P(z)=0$ by the Jacobi equation.
It is nonnegative when $a\in G_x$.  When $a\in I(x)$, it is
strictly positive: equality would make the nonzero field $J_a$ a
Jacobi field with zero endpoints, contrary to nonconjugacy.
At $a=z$, integration by parts gives
\[
 Q_z(E_z,J_z)
 =[\langle E_z,D_\tau J_z\rangle]_0^1
 =-\langle\xi,k\rangle_x.
\]
For $a\in I(x)$ near $z$, \cref{lem:hessian-index} can be tested
on $E_a+\lambda J_a$.  Minimizing the resulting quadratic polynomial
in $\lambda\in\R$ gives
\begin{equation}
 \begin{split}
 A_x(a)(\xi,\xi)
 &\leq Q_a(E_a,E_a)
       -\frac{Q_a(E_a,J_a)^2}{P(a)}\\
 &\leq C_0-\frac{c_0}{P(a)},
 \qquad c_0>0.
 \end{split}
 \label{eq:trial-hessian}
\end{equation}
The last inequality holds in one fixed neighborhood: smoothness
bounds $Q_a(E_a,E_a)$ from above and bounds $Q_a(E_a,J_a)^2$
away from zero there.

\medskip\noindent
\emph{Reduction to a one-dimensional kernel.}
For each fixed $\xi\perp e$, \cref{lem:transverse-concavity} on the
contracted segment gives
\[
 A_x(rz)(\xi,\xi)
 \geq \frac{l_1}{l_0+l_1}A_x(rb_0)(\xi,\xi)
     +\frac{l_0}{l_0+l_1}A_x(rb_1)(\xi,\xi).
\]
The two endpoint velocities are nonconjugate.  Smoothness of the
branch Hessians therefore makes the right-hand side bounded from
below as $r\uparrow1$.  On the other hand, $P(rz)>0$ and
$P(rz)\to P(z)=0$.  Inequality \eqref{eq:trial-hessian} would
contradict this lower bound if $\langle\xi,k\rangle_x\ne0$.
It follows that every $k\in\ker d_z\exp_x$ is orthogonal to
$e^\perp$, and hence belongs to $\Span\{e\}$.
The kernel is nontrivial, so it is exactly $\Span\{e\}$.
In particular, $e$ is a kernel vector and $z\perp e$.

\medskip\noindent
\emph{A joint tangential and radial estimate.}
Apply the trial-field construction with $k=\xi=e$, and retain the
notation $P$ for this choice.  Since $z+te\in G_x$ for small
$t$ of both signs, the smooth function $t\mapsto P(z+te)$ has a
local minimum of value zero at $t=0$.  Its first derivative vanishes,
so, for some $\delta>0$ and $C>0$,
\[
 0\leq P(z+te)\leq Ct^2\qquad(|t|\leq\delta).
\]
Shrink $\delta$ and choose $r_0<1$ so that both $z+te$ and
$r(z+te)$ lie in the fixed neighborhood used in
\eqref{eq:trial-hessian} whenever $|t|\leq\delta$ and
$r_0<r<1$.  Choose a Lipschitz constant $L$ for $P$ on a slightly
larger ball and a bound $K$ for $|z+te|_x$ in this range.
Since $r(z+te)\in I(x)$, we obtain the simultaneous estimate
\begin{equation}
 \begin{split}
 0<P\bigl(r(z+te)\bigr)
 &\leq P(z+te)+L(1-r)|z+te|_x\\
 &\leq Ct^2+LK(1-r)
 \leq C'\bigl(t^2+1-r\bigr).
 \end{split}
 \label{eq:quadratic-index}
\end{equation}
Combining this with \eqref{eq:trial-hessian}, and adjusting positive
constants, gives
\begin{equation}
 A_x\bigl(r(z+te)\bigr)(e,e)
 \leq C_1-\frac{c_1}{t^2+1-r},\qquad c_1>0,
 \label{eq:singular-upper}
\end{equation}
for all these $t,r$.  In particular, the constants do not depend
on the manner in which the two parameters approach $0$ and $1$.

\medskip\noindent
\emph{The radial identity gives a contradictory lower bound.}
Put $B_x(a)=g_x-A_x(a)$ and, for $0<r<1$, define
\[
 f_r(t)=B_x\bigl(r(z+te)\bigr)(z,z),
 \qquad -l_0\leq t\leq l_1.
\]
Because $z\perp e$, \cref{lem:transverse-concavity} with test
vector $z$ and velocity direction $re$ makes $f_r$ convex.
By \eqref{eq:radial-hessian}, $f_r(0)=0$.
Nonconjugacy at $b_1=z+l_1e$ bounds $f_r(l_1)$ from above
uniformly as $r\uparrow1$.  The convexity inequality on $[0,l_1]$
therefore yields
\[
 f_r(t)\leq\frac{t}{l_1}f_r(l_1)\leq C_2t
 \qquad(0<t\leq l_1)
\]
for $r$ close to $1$, with $C_2\geq0$ independent of $r$.
The radial identity also says
\begin{equation}
 B_x(a)(a,\cdot)=0,
 \qquad
 f_r(t)=t^2B_x\bigl(r(z+te)\bigr)(e,e).
 \label{eq:radial-null}
\end{equation}
Indeed, for $a=r(z+te)$ and $r>0$, the first identity annihilates
$z+te$ in either slot.  Inserting $z=(z+te)-te$ in both slots
gives the second identity, with no factor of $r$ remaining.
As $|e|_x=1$, we conclude that
\begin{equation}
 A_x\bigl(r(z+te)\bigr)(e,e)\geq1-\frac{C_2}{t}
 \qquad(0<t\leq l_1).
 \label{eq:singular-lower}
\end{equation}
Choose $r=1-t^2$ for sufficiently small positive $t$.  Both
\eqref{eq:singular-upper} and \eqref{eq:singular-lower} apply and
would imply
\[
 1-\frac{C_2}{t}\leq C_1-\frac{c_1}{2t^2}.
\]
Multiplying by $t^2$ and letting $t\downarrow0$ is impossible,
since $c_1>0$.  This contradiction proves the proposition.
\end{proof}


\begin{corollary}[Nonconjugacy of a minimizing convex hull]
\label{geo:first-cut}
Let $a_1,\ldots,a_m\in T_xM$ be nonconjugate velocities, with $m\geq1$,
and suppose
\[
 K=\operatorname{co}\{a_1,\ldots,a_m\}\subset G_x.
\]
Then every velocity in $K$ is nonconjugate.
\end{corollary}

\begin{proof}
We induct on the number of positive weights in a convex representation.
A representation with one positive weight gives one of the prescribed
generators. For a representation
\[
 w=\sum_{i=1}^k\theta_i a_i,\qquad
 \theta_i>0,\qquad \sum_{i=1}^k\theta_i=1,\qquad k\geq2,
\]
write
\[
 w=(1-\theta_k)w'+\theta_k a_k,\qquad
 w'=\sum_{i=1}^{k-1}\frac{\theta_i}{1-\theta_k}a_i.
\]
The vector $w'$ is nonconjugate by induction, and $a_k$ is
nonconjugate by hypothesis. Their entire segment lies in $K\subset G_x$,
so \cref{prop:segment-nonconjugacy} makes $w$ nonconjugate as well.
Discarding zero weights gives the assertion for every point of $K$.
\end{proof}

In the continuation argument, membership of the entire hull in $G_x$
will follow by taking limits from smaller scales. The corollary then
removes conjugacy without assuming convexity of the full minimizing
domain.
